\documentclass[11pt]{article}
\usepackage[utf8]{inputenc}\usepackage[T1]{fontenc}\usepackage[spanish,es-tabla]{babel}
\usepackage[a3paper,landscape,margin=16mm,headheight=15pt,headsep=6mm,footskip=9mm]{geometry}
\usepackage[scaled]{helvet}\renewcommand{\familydefault}{\sfdefault}
\usepackage{graphicx,xcolor,tabularx,array,fancyhdr,hyperref}
\definecolor{azul}{RGB}{35,91,133}\hypersetup{colorlinks=true,urlcolor=azul,pdftitle={Limbani R03 - Planos académicos},pdfauthor={Aquiles Taylor Ramos Yapo}}
\pagestyle{fancy}\fancyhf{}\fancyhead[L]{\textbf{MINICENTRAL LIMBANI -- PROYECTO ACADÉMICO DE DISEÑO}}
\fancyhead[R]{R03 | 14.10.2026}\fancyfoot[L]{Aquiles Taylor Ramos Yapo | Estudiante 227368 | Centrales Hidroeléctricas | Grupo A}
\fancyfoot[R]{A3 horizontal | Lámina \thepage/8}
\setlength{\parindent}{0pt}\setlength{\parskip}{4pt}
\newcommand{\dato}[1]{\csname dato#1\endcsname}\input{generado/valores}
\newcommand{\titulo}[2]{{\LARGE\bfseries\color{azul}#1\quad #2}\par\medskip}
\newcommand{\nota}[1]{\par\small\textbf{Condición de uso:} #1\par}
\begin{document}
\titulo{PL-01}{Ubicación, cuenca y ámbito de suministro}
\begin{minipage}[t]{.57\linewidth}\centering\includegraphics[width=.94\linewidth,height=205mm,keepaspectratio]{figuras/cuenca-r03.png}\end{minipage}\hfill
\begin{minipage}[t]{.40\linewidth}\centering\includegraphics[width=.96\linewidth,height=205mm,keepaspectratio]{figuras/corredor-r03.png}\end{minipage}
\nota{Perú, Puno, Sandia, Limbani. WGS84 / UTM zona 19 sur, EPSG:32719; coordenadas en metros. Norte y escala gráfica en mapas. Fuentes: PISCO SENAMHI, Terrain Tiles/Mapzen y OpenStreetMap contributors (ODbL). Área de captación de gabinete \dato{Acapt} km\(^2\); referencia geodésica PISCO \dato{Aref} km\(^2\). Los enlaces entre localidades son relaciones funcionales, no un trazado replanteado de línea. Lince carece de ubicación puntual confirmada.}
\newpage
\titulo{PL-02}{Emplazamiento y perfil hidráulico propuesto}
\begin{minipage}[t]{.50\linewidth}\centering\includegraphics[width=.96\linewidth,height=180mm,keepaspectratio]{figuras/emplazamiento-r03.png}\end{minipage}\hfill
\begin{minipage}[t]{.47\linewidth}\centering\includegraphics[width=.98\linewidth,height=180mm,keepaspectratio]{figuras/perfil-hidraulico-r03.png}\par\raggedright
\textbf{Referencias UTM 19S:} CM-A 425580 E / 8436580 N (memoria de línea); CM-B 425715 E / 8436650 N (CAD). Captación propuesta 425542 E / 8436231 N, ajuste cartográfico aproximado.\par
\textbf{Geometría de estudio:} canal 150 m, tubería \dato{LongPipe} m, salto bruto \dato{Hb} m. Referencias SRTM90 de terreno: 3285 y 3239 m; consumo supuesto antes de cámara 0.7 m.\end{minipage}
\nota{Mapa: norte, metros y escala gráfica. Perfil: esquema funcional con cotas de gabinete; no rasante de campo. La cámara de carga y restitución no tienen coordenadas confirmadas. Las dos CM difieren cerca de 152 m. No se reemplazan por el centro del pueblo. Catastro original leído sin modificar; contiene 482 entidades Casas, no un padrón de contratos.}
\newpage
\titulo{PL-03}{Obras hidráulicas y dimensiones de estudio}
\centering\includegraphics[width=.96\linewidth,height=146mm,keepaspectratio]{figuras/obras-r03.png}\par\raggedright
\begin{tabularx}{\linewidth}{>{\bfseries}l X}
Captación & Lateral por evaluar; reja bruta de estudio \dato{RejaArea} m\(^2\). Nivel y sección del río pendientes; paso ecológico independiente.\\
Canal & 150 m, sección rectangular b=1.60 m, altura=1.30 m, S=0.002, n=0.015. Tirante \dato{CanalY} m a 1.90 m\(^3\)/s.\\
Desarenador & Dos cámaras de 29 por 4 m, tirante 1.80 m; cada cámara admite todo el caudal de diseño. Granulometría y purga por verificar.\\
Cámara & 8 por 5 m, profundidad operativa 2.5 m; 60 m\(^3\) de regulación supuesta. Alivio para caudal derivado, no para avenida del río.\\
Tubería / restitución & Di=1000 mm, L=350 m; descarga y niveles aguas abajo sin levantamiento. Caudal nominal \dato{Qnom} m\(^3\)/s; límite 1.90.\\
\end{tabularx}
\nota{Diagrama sin escala. Dimensiones en metros, salvo indicación. No se definen azud, sección de cauce ni cotas de obra listas para construir. Avenida, socavación, nivel de inundación y compatibilidad del canal con la ladera requieren estudio propio.}
\newpage
\titulo{PL-04}{Tubería, transitorios y criterios de soporte}
\begin{minipage}[t]{.61\linewidth}\centering\includegraphics[width=.97\linewidth,height=190mm,keepaspectratio]{figuras/transitorio-r03.png}\end{minipage}\hfill
\begin{minipage}[t]{.36\linewidth}\raggedright
\large\textbf{Especificación académica}\par\normalsize
Acero soldado; Di nominal 1000 mm, De 1032 mm, espesor nominal 16 mm. Pared mínima nueva 14 mm y final 13 mm; tolerancia y corrosión supuestas.\par
Separación preliminar de silletas 4 m; 89 apoyos para metrado. Flexión local y fundaciones pendientes.\par
Presión estática aproximada 4.44 bar. Cierre de 8 s: máxima \dato{PNslow} bar en perfil ideal. Primer pico brusco válido: \dato{PNfast} bar.\par
Válvulas/accesorios PN40 candidatos. No se certifica separación de columna, vacío ni colapso de cavidades.\par
ELC absorbe rechazo mientras el regulador reduce caudal. Cierre mecánico controlado independiente.\par
Expansión libre de estudio \dato{Expansion} mm para 40 grados. Anclajes y juntas coordinados con movimiento.\end{minipage}
\nota{Gráficos calculados, no detalle constructivo de soporte. MOC monofásico detenido al superar el límite de presión baja. Los anclajes a roca requieren EMS, ensayos y dimensionamiento. Caudal de comprobación 1.90 m\(^3\)/s; perfil real y curva de cierre pendientes.}
\newpage
\titulo{PL-05}{Casa de máquinas y patio de equipos}
\centering\includegraphics[width=.92\linewidth,height=205mm,keepaspectratio]{figuras/casa-maquinas-r03.png}\par\raggedright
\nota{Planta dimensional de estudio en metros; escala mediante ejes del dibujo. Casa 20 por 12 m, 240 m\(^2\). Tres grupos hidroeléctricos de 200 kW, generadores 250 kVA/380 V/900 rpm. Patio de T1 de 800 kVA y tres módulos de respaldo de 200 kW. Frente de tableros reservado 1.5 m; verificar condiciones CNE y equipos reales. Posición, orientación geográfica, cimentación, drenaje y nivel de inundación sin levantamiento. Ventilación requerida \dato{AirFlow} m\(^3\)/s, por verificar con curvas.}
\newpage
\titulo{PL-06}{Diagrama unifilar eléctrico}
\centering\includegraphics[width=.96\linewidth,height=188mm,keepaspectratio]{figuras/unifilar-r03.png}\par\raggedright
\nota{Esquema sin escala. Frecuencia 60 Hz, operación inicial aislada. Barra: potencia combinada hidro/diésel limitada a 600 kW. T1: 800 kVA, 10/0.38 kV, AT delta / BT estrella; desplazamiento horario por especificar. Red primaria 8 km AAAC 95 mm\(^2\); seis SED total 730 kVA. TSA 20 kVA, 380/380--220 V, neutro de utilización propio. Interruptores candidatos: 630 A por hidro y 1600 A de barra, Icu 25 kA. Cortocircuito inicial máximo supuesto \dato{IKmax} kA; selección final y ajustes requieren datos de fabricante. Sin punto de conexión SEIN aprobado.}
\newpage
\titulo{PL-07}{Regulación y funciones de protección}
\centering\includegraphics[width=.94\linewidth,height=150mm,keepaspectratio]{figuras/regulacion-r03.png}\par\raggedright
\begin{tabularx}{\linewidth}{>{\bfseries}l X}
Velocidad / frecuencia & Un líder coordinado y reparto entre grupos; ELC de actuación rápida y gobierno hidráulico con cierre lento. Balastos 3 por 200 kW, capacidad térmica transitoria por verificar.\\
Tensión / reactivos & AVR y reparto de reactivos. Sincronismo 25 para acoplar grupos a la barra interna, sin conexión inicial al SEIN.\\
Generadores & 87G, 51V, 64G/59N, 32, 40, 46, 27/59, 81 y sobrevelocidad independiente. Ajustes y funciones definitivas según equipos.\\
Transformación / red & 87T y coordinación de sobrecorriente/tierra según impedancia de neutro, decremento de fuentes y falla remota. No basta 50/51 simple.\\
\end{tabularx}
\nota{Diagrama sin escala; nombres ANSI describen funciones, no ajustes aprobados. El control no reemplaza protecciones independientes. Se requieren ensayos de rechazo, escalones de carga, falla de regulador, transferencia y operación con respaldo. Batería 24 V/100 Ah para control; no sostiene motores auxiliares.}
\newpage
\titulo{PL-08}{Puesta a tierra y funciones de neutro}
\centering\includegraphics[width=.94\linewidth,height=197mm,keepaspectratio]{figuras/puesta-tierra-r03.png}\par\raggedright
\nota{Planta conceptual con dimensiones en metros y ejes de escala. Malla 20 por 12 m, paso 4 m, 152 m horizontales, ocho electrodos de 2.4 m; cobre 50 mm\(^2\) propuesto sin verificación térmica definitiva. Neutro de generación por impedancia coordinada; 10 kV necesita dispositivo de neutro si se adopta esa puesta a tierra; utilización BT con neutro y PE definidos por separado. Se requieren resistividad estratificada, corriente a tierra y tiempo de despeje. GPR no es tensión de contacto; una resistencia global de 5 ohm no demuestra seguridad. CNE Suministro 036: verificar paso, contacto y transferencia.}
\end{document}
